\section{中兴通讯：经营证据已闭环，但赔率未过关}

现价40.53元，熊／基准／牛情景为25.00/37.00/51.90元，35\%／45\%／20\%概率值35.78元，相对现价-11.7\%。2025年算力相关收入同比约+150\%、占营收24.6\%，服务器及存储收入+200\%以上、数据中心产品+50\%；2026Q1算力占比升至27\%。国金原始PDF明确把头部互联网、运营商、政务与金融智算项目列为需求来源，并给出2026E收入1728.08亿元、净利73.73亿元、EPS 1.541元、经营现金流／股1.80元与毛利率31.3\%。具名终端客户、设备量、合同ASP和算力分部毛利在已检查季报及两份原始PDF中没有单独披露，所以不使用算力分部高倍数；旧48.75元目标撤销，当前只用合并EPS情景，且概率值低于现价。

\section{江波龙：上半年兑现不等于周期永久化}

江波龙H1归母净利中值101亿元、扣非97.5亿元，Q1归母38.62亿元，隐含Q262.38亿元，主营兑现强；但Q1经营现金流为-28.75亿元，一年位置约79\%，当前587.60元已包含相当多周期持续性假设。

\begin{exhibitbox}[Exhibit 14：江波龙周期敏感度]
\begin{tabularx}{\textwidth}{L{1.6cm}R{1.5cm}R{1.4cm}R{1.3cm}R{1.4cm}X}
\toprule
\textbf{情景} & \textbf{全年净利} & \textbf{EPS} & \textbf{PE} & \textbf{价值} & \textbf{含义} \\
\midrule
熊 & 131亿元 & 约31元 & 11x & 341元 & 库存红利消退、价格回落 \\
基准 & 161亿元 & 约38元 & 14x & 532元 & H2仍盈利但不继续加速 \\
牛 & 191亿元 & 约45元 & 17x & 765元 & 合同价、供给与订单延续 \\
\bottomrule
\end{tabularx}
\end{exhibitbox}

概率值511.75元，相对现价-12.9\%；旧802.30元目标撤销。官方预告明确披露与多家全球存储晶圆原厂续签LTA／MOU、自研SPU／HLC、自有高端封测以及与AMD联合调优，后者令端侧AI产品DRAM使用量下降约40\%。国信与爱建原始PDF补充企业级存储+93.3\%、平台认证、通信／金融／互联网客户链和NAND／DRAM价格指数。合同价、库存成本层和H2具名订单额属于正式披露边界，因此模型使用11x／14x／17x周期PE并保留35\%熊市概率；升级要求H2扣非、现金流和存货减值共同验证，而不是等待一个可能永不公开的客户名称。

\section{陕西煤业：行业静默成立，个股赔率未过关}

陕西煤业H1归母／扣非净利中值114.58／97.46亿元，行业阶段为静默吸纳，但现价22.92元高于概率值19.76元。熊／基准／牛为13.60/20.80/26.40元。股息提供下行支撑，但只有煤价、2026E EPS和分红持续性同时高于当前券商基准，才可从行业验证升级为正式配置。

\begin{riskbox}[降级原则]
条件观察并非“没有证据”：中兴和江波龙都已完成直接披露、代理证据和正式边界记录。但只要概率值低于现价，或现金流／库存／利润门槛尚未兑现，就不能用牛市目标替代正式目标。
\end{riskbox}
